%% Full exponential, classifier, and power-object constructions.

\section{The elementary topos structure of a presheaf category}

\begin{annexitem}{annex:p10.psh.3}{thm:p10-presheaf-exponentials}{Evaluation and the two transpose identities}
\begin{proof}
\paragraph{Sets of components and the restriction action.}
Since \(\mathcal C\) is small, the set
\(\operatorname{Nat}(\mathrm y c\times X,Y)\) is a subset of
\[
       \prod_{d\in\operatorname{Ob}\mathcal C}
          \mathsf{Set}(\mathcal C(d,c)\times X(d),Y(d)).
\]
Every factor and the indexing set belong to the fixed universe;
the product and its naturality-defined subset do too.
This establishes the required component size.

For \(u:d\to c\), the map
\((\mathrm y u)_e(v)=u\circ v\) is natural in \(e\), by associativity
of composition.  Thus precomposing a natural \(\alpha\) with
\(\mathrm y u\times1_X\) produces a natural transformation.
The identity action fixes \(\alpha\).
For \(v:e\to d\), composition gives
\[
             (Y^X)(v)(Y^X)(u)(\alpha)=(Y^X)(u\circ v)(\alpha).
\]
Hence \(Y^X\) is a presheaf.

\paragraph{Evaluation is natural.}
For \(u:d\to c\), \(\alpha\in(Y^X)(c)\), and \(x\in X(c)\),
naturality of \(\alpha:\mathrm y c\times X\to Y\) gives
\[
\begin{aligned}
 Y(u)(\operatorname{ev}_c(\alpha,x))
   &=Y(u)(\alpha_c(1_c,x))\\
   &=\alpha_d(u,X(u)(x))\\
   &=\operatorname{ev}_d((Y^X)(u)(\alpha),X(u)(x)).
\end{aligned}
\]
This is exactly the evaluation naturality square.

\paragraph{Both naturalities of the transpose.}
Fix \(h:Z\times X\to Y\), \(c\), and \(z\in Z(c)\).
The proposed component at \(d\) is a function
\(\mathcal C(d,c)\times X(d)\to Y(d)\).
For \(v:e\to d\), naturality of \(h\) and functoriality of \(Z\)
give
\[
\begin{aligned}
 Y(v)(h_d(Z(u)(z),x))
   &=h_e(Z(v)Z(u)(z),X(v)(x))\\
   &=h_e(Z(u\circ v)(z),X(v)(x)).
\end{aligned}
\]
This proves that \((\Lambda h)_c(z)\) is a natural transformation
\(\mathrm y c\times X\to Y\).

For \(u:d\to c\), the equality
\[
                 (Y^X)(u)((\Lambda h)_c(z))
                           =(\Lambda h)_d(Z(u)(z))
\]
is checked at \(v:e\to d\) and \(x\in X(e)\).
Both sides give \(h_e(Z(u\circ v)(z),x)\).
Thus \(\Lambda h\) is natural in the base object too.

\paragraph{The inverse identities and parameter naturality.}
Evaluation at \(1_c\) gives
\[
     \operatorname{ev}_c((\Lambda h)_c(z),x)
                     =h_c(Z(1_c)(z),x)=h_c(z,x).
\]
This is the first inverse identity.
Conversely let \(k:Z\to Y^X\), and put
\(h=\operatorname{ev}\circ(k\times1_X)\).
For \(u:d\to c\), naturality of \(k\) gives
\(k_d(Z(u)(z))=(Y^X)(u)(k_c(z))\).
Consequently
\[
\begin{aligned}
 [(\Lambda h)_c(z)]_d(u,x)
   &=h_d(Z(u)(z),x)\\
   &=[k_d(Z(u)(z))]_d(1_d,x)\\
   &=[(Y^X)(u)(k_c(z))]_d(1_d,x)\\
   &=[k_c(z)]_d(u,x).
\end{aligned}
\]
All components agree, proving the second inverse identity and
the uniqueness of the transpose.

For \(a:Z'\to Z\), at \(z'\in Z'(c)\), \(u:d\to c\), and
\(x\in X(d)\), the two transpose formulas give respectively
\(h_d(a_d(Z'(u)(z')),x)\) and
\(h_d(Z(u)(a_c(z')),x)\).  Naturality of \(a\) equates them,
so
\[
            \Lambda(h\circ(a\times1_X))=(\Lambda h)\circ a.
\]
For \(b:Y\to Y'\), define \(b^X:Y^X\to(Y')^X\) by
\((b^X)_c(\alpha)=b\circ\alpha\).
It is natural because this postcomposition commutes with the
restriction precompositions.  It preserves identity and composition
in \(b\), and
\[
                            \Lambda(b\circ h)=b^X\circ\Lambda h.
\]
These equations prove naturality of the hom-set bijection in
both indicated variables.  With the already proved pointwise
terminal object and products, they establish cartesian closure.

If \(X=0\), each \(\mathrm y c\times X\) is initial, so each
exponential component is a singleton and \(Y^0\) is terminal.
If \(Y=0\), a component set consists of natural maps into the
empty presheaf: it is empty if any domain component is nonempty,
and otherwise contains the unique map.  The proof above only
quantifies over existing elements and maps, so covers both
possibilities without selecting an element of an empty set.
If \(\mathcal C\) is empty, there is one presheaf and one
natural transformation; all the displayed families have no
components and both universal-property hom-sets are singletons.
\end{proof}
\end{annexitem}

\begin{annexitem}{annex:p10.psh.4}{thm:p10-presheaf-topos}{The unique sieve characteristic map}
\begin{proof}
\paragraph{The truth-value presheaf.}
The arrows into \(c\) form a small subset of the arrow set of
\(\mathcal C\), so its collection of sieves is small.
If \(v:e\to d\) belongs to \(\Omega(u)(S)\) and \(w:b\to e\),
then \(u\circ v\in S\) implies \(u\circ v\circ w\in S\).
Thus \(\Omega(u)(S)\) is a sieve.
The identity restriction fixes a sieve, and
\[
                      \Omega(v)\Omega(u)(S)=\Omega(u\circ v)(S)
\]
follows by testing membership of each arrow into the new domain.
These are the presheaf laws.

A sieve on \(c\) containing \(1_c\) contains
\(1_c\circ u=u\) for every arrow into \(c\), so is maximal.
Every restriction sends the maximal sieve to the maximal sieve.
Hence \(\mathsf{true}\) is natural.
Its singleton-domain components are injective, so it is monic by
Theorem~\ref{thm:p10-presheaf-mono-epi}.
The equivalence between sieves and subpresheaves of \(\mathrm y c\)
is immediate in both directions: restricting an arrow is
precomposition, which is exactly sieve closure.

\paragraph{Existence of the characteristic map.}
Let \(A\subseteq X\) be a subpresheaf.
If \(u:d\to c\) belongs to \((\chi_A)_c(x)\) and \(v:e\to d\),
then \(X(u)(x)\in A(d)\).  Restriction stability of \(A\) gives
\[
                      X(v)X(u)(x)=X(u\circ v)(x)\in A(e).
\]
This proves sieve closure.
For \(u:d\to c\), membership of \(v:e\to d\) in either
\(\Omega(u)((\chi_A)_c(x))\) or
\((\chi_A)_d(X(u)(x))\) says exactly
\(X(u\circ v)(x)\in A(e)\).
Thus \(\chi_A\) is natural.

Its value at \(x\in X(c)\) is maximal if and only if
\(x\in A(c)\).  The forward implication tests \(1_c\);
the reverse uses closure of \(A\) under every restriction.
Thus the component square with domain \(A(c)\) is a pullback
in \(\mathsf{Set}\): a function into \(X(c)\) whose composite
with \((\chi_A)_c\) is constantly \(t(c)\) has every value
in \(A(c)\), and factors through that subset uniquely.
This proves the universal property for the displayed domain,
without identifying it literally with a chosen set of pullback
pairs.  Pointwise limits, as in
Theorem~\ref{thm:p10-presheaf-limits}, prove that the displayed
square is a pullback in \(\widehat{\mathcal C}\).

\paragraph{Uniqueness and arbitrary monomorphisms.}
Suppose \(\chi:X\to\Omega\) also has \(A\) as that pullback.
For \(u:d\to c\) and \(x\in X(c)\), the defining restriction,
naturality, and the pullback description give
\[
\begin{aligned}
 u\in\chi_c(x)
 &\Longleftrightarrow 1_d\in\Omega(u)(\chi_c(x))\\
 &\Longleftrightarrow 1_d\in\chi_d(X(u)(x))\\
 &\Longleftrightarrow \chi_d(X(u)(x))=t(d)\\
 &\Longleftrightarrow X(u)(x)\in A(d).
\end{aligned}
\]
Thus \(\chi=\chi_A\).

For a general mono \(m:B\to X\), use the canonical image
\(A(c)=m_c[B(c)]\).  The componentwise mono criterion makes
the induced map \(B\to A\) a componentwise bijection; it is
a natural isomorphism by
Theorem~\ref{thm:p10-presheaf-mono-epi}.
Composing the inclusion \(A\to X\) with this isomorphism
therefore yields the asserted pullback square for \(m\).
Conversely any characteristic square for \(m\), computed
pointwise, has image
\(\{x:\chi_c(x)=t(c)\}\) in each component.  Hence its image
is \(A\), and the just-proved uniqueness applies.
Isomorphic monos over \(X\) have the same image, so this
classification is independent of the representative.

For later use, if \(a:Z\to X\), the inverse-image
subpresheaf \(a^*A\) has characteristic map
\[
                               \chi_{a^*A}=\chi_A\circ a.
\]
Indeed, for \(z\in Z(c)\), membership of \(u:d\to c\) on
the left says \(a_d(Z(u)(z))\in A(d)\).  Naturality of \(a\)
identifies this with \(X(u)(a_c(z))\in A(d)\), the condition
on the right.  This proves the naturality of classification
under arbitrary presheaf maps.

The preceding proof supplies a subobject classifier.
Finite limits are already available pointwise, and the
previous theorem supplies exponentials, so the stated
elementary-topos definition is satisfied.
If the base is empty, all component assertions are vacuous.
There is only the empty functor and its identity map; that
map is the truth map and classifies the sole subobject.
Thus the degenerate case satisfies the same universal properties.
\end{proof}
\end{annexitem}

\begin{annexitem}{annex:p10.psh.5}{thm:p10-presheaf-power-objects}{Power-object classification and membership}
\begin{proof}
\paragraph{Small subobject sets and the universal correspondence.}
Using canonical image representatives, a subobject of a
presheaf \(W\) is a restriction-stable family of subsets
\(A(c)\subseteq W(c)\).  Such families form a subset of
\[
                    \prod_{c\in\operatorname{Ob}\mathcal C}
                                   \mathcal P(W(c)),
\]
which is a small set.  This proves the required size of
\(\operatorname{Sub}(W)\), without choosing representatives
from a collection of unrelated mono domains.

For \(k:Z\to\Omega^X\), uncurrying gives
\(h=\operatorname{ev}\circ(k\times1_X):Z\times X\to\Omega\).
The classifier associates to \(h\) exactly its inverse-image
subobject \(h^*(\mathsf{true})\).
Since \(\in_X=\operatorname{ev}^*(\mathsf{true})\), the
componentwise inverse-image definition gives
\[
                  h^*(\mathsf{true})=(k\times1_X)^*(\in_X).
\]
Conversely, for a subobject \(S\) of \(Z\times X\), take its
unique characteristic map \(\chi_S\), and transpose it.
The classifier and exponential bijections are both inverse
bijections, so these procedures are inverse and give
uniqueness of the resulting \(k\).

For \(a:Z'\to Z\), the classifier naturality just proved and
the transpose naturality show that \(\Lambda(\chi_S)\circ a\)
classifies \((a\times1_X)^*S\).
Equivalently,
\[
 ((k\circ a)\times1_X)^*(\in_X)
                =(a\times1_X)^*((k\times1_X)^*(\in_X)).
\]
This proves the asserted naturality in \(Z\) and the full
power-object universal property.

\paragraph{Components and restrictions.}
By the actual exponential construction,
\[
             P(X)(c)=\operatorname{Nat}(\mathrm y c\times X,\Omega).
\]
Apply the classifier bijection to the presheaf
\(\mathrm y c\times X\).  It sends each such natural
transformation to its inverse-image subpresheaf, and
provides the stated canonical identification with
\(\operatorname{Sub}(\mathrm y c\times X)\).
The restriction along \(u:d\to c\) is precomposition
with \(\mathrm y u\times1_X\).  Naturality of classification
therefore identifies it with pullback along that map.
At \(e\), the map is
\((v,x)\mapsto(u\circ v,x)\), proving the explicit formula
for \((S|u)(e)\).

\paragraph{Membership and the classifying map of a relation.}
Let \(\alpha:\mathrm y c\times X\to\Omega\) correspond to
the subpresheaf \(S\).
By evaluation,
\[
\begin{aligned}
 (\alpha,x)\in(\in_X)(c)
 &\Longleftrightarrow \alpha_c(1_c,x)=t(c)\\
 &\Longleftrightarrow (1_c,x)\in S(c).
\end{aligned}
\]
This proves the membership formula under the declared
component identification.
Restriction stability can also be seen directly:
if \((1_c,x)\in S(c)\) and \(u:d\to c\), restricting in
\(\mathrm y c\times X\) gives \((u,X(u)(x))\in S(d)\).
The formula for \(S|u\) then gives
\((1_d,X(u)(x))\in(S|u)(d)\), exactly the required
membership after restriction.

For \(S\subseteq Z\times X\) and \(z\in Z(c)\), the natural
map \(\bar z:\mathrm y c\to Z\) has components
\(\bar z_d(u)=Z(u)(z)\).
Its naturality follows from functoriality of \(Z\).
The pullback \(S_z=(\bar z\times1_X)^*S\) has precisely
the displayed components, so is a subpresheaf.
By classifier naturality, its characteristic map is
\(\chi_S\circ(\bar z\times1_X)\).
The transpose formula shows that this is the element
\((\Lambda\chi_S)_c(z)\) of \(P(X)(c)\).
Hence the stated description is exactly the unique
classifying map, not an additional choice.

In particular, pulling membership back along this map
recovers \(S\): at \(c\), a pair \((z,x)\) satisfies the
pulled-back membership condition if and only if
\((1_c,x)\in S_z(c)\), which is equivalent to
\((Z(1_c)(z),x)=(z,x)\in S(c)\).
These are equality of the represented subobjects; for
an arbitrary mono the canonical image convention gives
the same subobject classification.

When \(X=0\), \(P(X)=\Omega^0\) is terminal and membership
has initial domain.  Each \(Z\times0\) has only its empty
subpresheaf, so both sides of the classification bijection
are singletons.  If the base is empty there is just one
presheaf and one subobject, and all statements again
reduce to singleton bijections.  The constructions use
only natural component functions, images, inverse images,
and universal properties; no family of witnesses or
mono representatives is selected by the Axiom of Choice.
\end{proof}
\end{annexitem}
